\section{Provenance and contribution disclosure}\label{app:provenance}

This companion was separated from the unified Wilson--Loewner
manuscript on 20 September 2026 by OpenAI GPT-6 (Codex), for
Edward Baker. Version 0.5 shares the exposition's version number
to identify the matched pair. The configured reasoning effort
was not exposed to the assistant in this session and is not
inferred. Both documents are working research drafts without
independent mathematical review.

Sections~S1--S8 preserve the detailed arithmetic, free-field,
contour, supersymmetry, endpoint and reflection calculations from
versions 0.1--0.2. The three principal local source notes are
\emph{Smooth Wilson variation and the constant Loewner driver},
\emph{Defect-compatible contour freedom and the endpoint pairing},
and \emph{Physical reflection, the reference junction, and a first
bulk response}, all dated 19 September 2026. The parent
Wilson-lines work supplies the angular regulator and auxiliary
Gaussian model; the Loewner and fractional-dimension investigations
supply the transfer factorization and scoped realization restrictions.
Section~S9 transcribes the inherited cumulative continuation
algebra~\cite{CumulativePath} and now reports the separately
certified fixed-window anchor and obstruction to the scalar append
estimate \cite{EMAAnchor,CumulativeAppend}. The subsequent internal operator-coupling bound is cited separately
\cite{CumulativeCoupling}; its proof stays in the critical-path investigation. The full program and hypotheses are in the
exposition~\cite{Exposition}.

The original assembly checked the printed conventions for Loewner
normalization, tangent and D3--D5 projectors, Baker endpoint
transformations and semicircle, Dorn--Pershin reflection, and the
Weil-form weights against the cited primary papers. This revision
does not claim a new exhaustive literature or priority survey.
No novelty is claimed for standard transport variation,
tangent-coupled supersymmetry, the ray endpoint selection rule,
or the classical Weil criterion.

Version 0.4 incorporates the fixed-window Wilson-response and
localization notes of 20 September 2026. Sections~S10--S11 add the
full contact/pole comparison, the defined Gaussian readout and its
scoped obstructions, and the auxiliary boson--fermion product.
The localization direction follows Edward Baker's proposal and
his earlier open-line research. Its physical mode pairing and
operator interpretation remain unproved. The additional literature
checks concern Pestun's fluctuation complex, Wang's defect bilocal
and determinant assumption, and the gamma-function identities.
The 31 additional diagnostics bring the total to 299; this revision
has not received independent mathematical review.

\nolinkurl{MANUSCRIPT_SOURCES.json} records assembly and revision
inputs. \nolinkurl{BUILD_RECORD.json} binds both documents and
supporting checks; each dated archive has a
\nolinkurl{SNAPSHOT.json}. These are identity records, not proof
certificates. The historical unified drafts remain unchanged.
Every TeX input needed to build either current document is in this
package; no shared-background or parent-investigation TeX file is
loaded at build time.

Version 0.5 adds Section~\ref{sec:localized-response}, based on the
predictive hemisphere, spatial radial-descent and Hodge-channel
notes of 20 September 2026. The localization direction follows
Edward Baker's suggestion. The free boundary and Hodge inputs come
from \cite{DPYHiggs,DedushenkoGluing}; the derived spatial restrictions
do not exclude an enlarged or interacting sector. The three
additional standard-library controls retain separate record formats
and are not folded into the historical 299-case count.
All new exposition and checks were prepared with OpenAI GPT-6
(Codex) assistance; the effort setting was not exposed and is not
inferred. No specialist mathematical or physical review is claimed.


The live revision of 21 September 2026 incorporates Baker's clarification
that a specified Loewner evolution generates the curve defining the quantum
observable, and that the immediate aim is its evolution equation in a fixed
theory. The new chord-completed Yang--Mills calculation, its ordered hierarchy
and quantitative remainder come from the research note \cite{GrowingTraceNote}.
They were derived and integrated with OpenAI GPT-6 (Codex) assistance; the
reasoning effort was not exposed. The field-theory regularity assumptions,
finite matrix diagnostics and outstanding specialist review are stated
separately. The archived version 0.5 and all earlier snapshots remain unchanged.
At Baker's request this is an unarchived live revision: no new dated draft or
version number has been saved. Arithmetic manuscript separation is deferred.
